\section*{Working note: explicit sextic partition with visible harmonic scaling}

The backend implementation constructs the standard sextic Cartesian components
directly in the principal-axis frame. Throughout this note:

\begin{itemize}
\item $i,j,k,\ldots$ denote normal modes;
\item $a,b,c$ denote principal-axis labels.
\end{itemize}

The reduced cubic force tensor used in the code is the fully symmetrized
Gaussian normal-mode quantity $\Phi^{(3),\mathrm{red}}_{ijk}$, stored over
ordered triples. Its partition is therefore:

\begin{itemize}
\item semi-diagonal sector: all triples with at least two equal indices;
\item genuine three-index sector: all triples with $i,j,k$ distinct.
\end{itemize}

\subsection*{Mode-resolved first-level object}

Define the backend-faithful first-level quantity
\begin{equation}
\widetilde{\mu}^{(i)}_{ab}
=
-\frac{(I_0\,\mu^{(i)}\,I_0)_{ab}}
{2\,M_a M_b\,\kappa_i},
\qquad
\kappa_i
=
\mathrm{DIDQ}_0\,(\mathrm{FACTG}\,\omega_i)^{3/2},
\end{equation}
where $I_0$ is the equilibrium inertia tensor, $M_a$ are the equilibrium
principal moments, and $\mu^{(i)}$ is the first normal-mode derivative of
$I^{-1}$.

To expose the harmonic power counting, introduce the rescaled quantity
\begin{equation}
\mathcal{M}^{(i)}_{ab}
=
\omega_i^{3/2}\widetilde{\mu}^{(i)}_{ab}.
\end{equation}

The quartic tensor entering the sextic formulas becomes
\begin{equation}
\tau_{abcd}
=
-2\sum_i
\frac{\mathcal{M}^{(i)}_{ab}\mathcal{M}^{(i)}_{cd}}{\omega_i^2}.
\label{eq:sextic_working_tau_mu1}
\end{equation}
Hence each quartic insertion is explicitly of order $\omega^{-2}$, so every
purely geometric sextic contribution built quadratically from $\tau$ is of
order $\omega^{-4}$.

\subsection*{Coriolis-dressed first-level object}

The quantity appearing in the sextic Coriolis geometry terms can be written as
\begin{equation}
\widetilde{\nu}^{(i)}_{abc}
=
\sum_j
\frac{2}{3}\,
\frac{\omega_i^2+2\omega_j^2}{\omega_i^{5/2}\omega_j^2}
\Big[
B_a\,\zeta^{(a)}_{ij}\,\mathcal{M}^{(j)}_{bc}
+
B_b\,\zeta^{(b)}_{ij}\,\mathcal{M}^{(j)}_{ca}
+
B_c\,\zeta^{(c)}_{ij}\,\mathcal{M}^{(j)}_{ab}
\Big].
\label{eq:sextic_working_nu_mu1}
\end{equation}
The sextic terms depending only on $\tau$ and $\widetilde{\nu}$ therefore
belong to the geometric sector and remain homogeneous of order $\omega^{-4}$.

\subsection*{Diagonal component $\Phi_{aaa}$}

The backend-faithful decomposition is
\begin{equation}
\Phi_{aaa}=X_1-X_2+X_3+X_4,
\end{equation}
with
\begin{align}
X_1 &= \frac{3}{16}\sum_x\frac{\tau_{aaax}^2}{B_x},
&
X_2 &= \frac{1}{2}\sum_i \omega_i\left(\widetilde{\nu}^{(i)}_{aaa}\right)^2,
\\
X_3 &= \frac{1}{6}\sum_{i,j,k}
\frac{\Phi^{(3),\mathrm{red}}_{ijk}}{\omega_i\omega_j\omega_k}\,
\mathcal{M}^{(i)}_{aa}\mathcal{M}^{(j)}_{aa}\mathcal{M}^{(k)}_{aa},
&
X_4 &= \frac{1}{4}\sum_{x\neq a}
\frac{\tau_{xaaa}^2}{B_a-B_x}.
\end{align}

Thus:
\begin{itemize}
\item $X_1$, $X_2$, and $X_4$ belong to the geometric sector and scale as
$\omega^{-4}$;
\item $X_3$ belongs to the cubic sector and scales as $\omega^{-3}$.
\end{itemize}

The cubic term splits as
\begin{equation}
X_3 = X_3^{(2)} + X_3^{(3)},
\end{equation}
with
\begin{align}
X_3^{(2)}
&=
\frac{1}{6}\sum_i
\frac{\Phi^{(3),\mathrm{red}}_{iii}}{\omega_i^3}
\left(\mathcal{M}^{(i)}_{aa}\right)^3
+\frac{1}{2}\sum_{i\neq j}
\frac{\Phi^{(3),\mathrm{red}}_{iij}}{\omega_i^2\omega_j}
\left(\mathcal{M}^{(i)}_{aa}\right)^2\mathcal{M}^{(j)}_{aa},
\\
X_3^{(3)}
&=
\sum_{i<j<k}
\frac{\Phi^{(3),\mathrm{red}}_{ijk}}{\omega_i\omega_j\omega_k}
\mathcal{M}^{(i)}_{aa}\mathcal{M}^{(j)}_{aa}\mathcal{M}^{(k)}_{aa}.
\end{align}

\subsection*{Mixed component $\Phi_{aab}$}

For $a\neq b$ the backend uses
\begin{equation}
\Phi_{aab}=Y_1-Y_2+Y_3+Y_4+Y_5+Y_6,
\end{equation}
where
\begin{align}
Y_1 &=
\frac{3}{32}\sum_x
\frac{T_{abx}^2+2\tau_{aaax}U_{abx}}{B_x},
&
T_{abx} &= \tau_{aabx}+2\tau_{abax},
&
U_{abx} &= \tau_{bbax}+2\tau_{abbx},
\\
Y_2 &=
\frac{3}{4}\sum_i \omega_i
\left[
3\left(\widetilde{\nu}^{(i)}_{aab}\right)^2
+2\,\widetilde{\nu}^{(i)}_{aaa}\widetilde{\nu}^{(i)}_{abb}
\right],
\\
Y_3 &=
\frac{1}{4}\sum_{i,j,k}
\frac{\Phi^{(3),\mathrm{red}}_{ijk}}{\omega_i\omega_j\omega_k}
\mathcal{M}^{(k)}_{aa}
\left[
\mathcal{M}^{(i)}_{aa}\mathcal{M}^{(j)}_{bb}
+4\mathcal{M}^{(i)}_{ab}\mathcal{M}^{(j)}_{ab}
\right],
\\
Y_4 &=
\frac{\tau_{aaab}\left(4\tau_{bbba}-3\tau_{aaab}\right)}
{8(B_a-B_b)},
\\
Y_5 &=
\sum_{x\neq a,b}
\frac{
\tau_{aaax}
\Big[
(B_a-B_x)(\tau_{bbax}+2\tau_{babx})
+(B_a-B_b)(\tau_{xxxa}-\tau_{aaax})
\Big]}
{4(B_a-B_x)^2},
\\
Y_6 &=
\sum_{x\neq a,b}
\frac{
(\tau_{aabx}+2\tau_{abax})
\Big[
(B_b-B_x)(\tau_{aabx}+2\tau_{abax})
+2(B_a-B_b)(\tau_{bbbx}-\tau_{xxxb})
\Big]}
{8(B_b-B_x)^2}.
\end{align}

Again:
\begin{itemize}
\item $Y_1$, $Y_2$, $Y_4$, $Y_5$, and $Y_6$ are geometric and scale as
$\omega^{-4}$;
\item $Y_3$ is cubic and scales as $\omega^{-3}$.
\end{itemize}

The cubic term splits as
\begin{equation}
Y_3 = Y_3^{(2)} + Y_3^{(3)},
\end{equation}
with
\begin{align}
Y_3^{(2)}
&=
\frac{1}{4}\sum_i
\frac{\Phi^{(3),\mathrm{red}}_{iii}}{\omega_i^3}
\mathcal{M}^{(i)}_{aa}
\left[
\mathcal{M}^{(i)}_{aa}\mathcal{M}^{(i)}_{bb}
+4\left(\mathcal{M}^{(i)}_{ab}\right)^2
\right]
\notag\\
&\quad
+\frac{1}{4}\sum_{i\neq j}
\frac{\Phi^{(3),\mathrm{red}}_{iij}}{\omega_i^2\omega_j}
\Big\{
\mathcal{M}^{(j)}_{aa}
\left[
\mathcal{M}^{(i)}_{aa}\mathcal{M}^{(i)}_{bb}
+4\left(\mathcal{M}^{(i)}_{ab}\right)^2
\right]
\notag\\
&\qquad\qquad
+\mathcal{M}^{(i)}_{aa}
\left[
\mathcal{M}^{(i)}_{aa}\mathcal{M}^{(j)}_{bb}
+4\mathcal{M}^{(i)}_{ab}\mathcal{M}^{(j)}_{ab}
\right]
\notag\\
&\qquad\qquad
+\mathcal{M}^{(i)}_{aa}
\left[
\mathcal{M}^{(j)}_{aa}\mathcal{M}^{(i)}_{bb}
+4\mathcal{M}^{(j)}_{ab}\mathcal{M}^{(i)}_{ab}
\right]
\Big\},
\\
Y_3^{(3)}
&=
\frac{1}{4}\sum_{i<j<k}
\frac{\Phi^{(3),\mathrm{red}}_{ijk}}{\omega_i\omega_j\omega_k}
\sum_{\pi\in S_3}
\mathcal{M}^{(\pi_3)}_{aa}
\left[
\mathcal{M}^{(\pi_1)}_{aa}\mathcal{M}^{(\pi_2)}_{bb}
+4\mathcal{M}^{(\pi_1)}_{ab}\mathcal{M}^{(\pi_2)}_{ab}
\right].
\end{align}

\subsection*{Fully mixed component $\Phi_{abc}$}

For three distinct axes the backend uses
\begin{equation}
\Phi_{abc}=Z_1-Z_2+Z_3+Z_4,
\end{equation}
with
\begin{align}
Z_1
&=
\frac{3}{16}\sum_x
\frac{
2V_x^2
+W_x^{(ab)}W_x^{(ac)}
+W_x^{(bc)}W_x^{(ba)}
+W_x^{(ca)}W_x^{(cb)}
}{B_x},
\\
V_x &= \tau_{xabc}+\tau_{xbca}+\tau_{xcab},
\\
W_x^{(ab)} &= \tau_{xabb}+2\tau_{xbab},
\qquad
W_x^{(ac)} = \tau_{xacc}+2\tau_{xcac},
\\
W_x^{(bc)} &= \tau_{xbcc}+2\tau_{xcbc},
\qquad
W_x^{(ba)} = \tau_{xbaa}+2\tau_{xaba},
\\
W_x^{(ca)} &= \tau_{xcaa}+2\tau_{xaca},
\qquad
W_x^{(cb)} = \tau_{xcbb}+2\tau_{xbcb},
\\
Z_2
&=
\frac{9}{2}\sum_i\omega_i
\Big[
2\left(\widetilde{\nu}^{(i)}_{abc}\right)^2
+\widetilde{\nu}^{(i)}_{bba}\widetilde{\nu}^{(i)}_{cca}
+\widetilde{\nu}^{(i)}_{ccb}\widetilde{\nu}^{(i)}_{aab}
+\widetilde{\nu}^{(i)}_{aac}\widetilde{\nu}^{(i)}_{bbc}
\Big],
\\
Z_3
&=
\frac{1}{2}\sum_{i,j,k}
\frac{\Phi^{(3),\mathrm{red}}_{ijk}}{\omega_i\omega_j\omega_k}
\Big[
\mathcal{M}^{(i)}_{aa}\mathcal{M}^{(j)}_{bb}\mathcal{M}^{(k)}_{cc}
+2\mathcal{M}^{(i)}_{aa}\mathcal{M}^{(j)}_{bc}\mathcal{M}^{(k)}_{bc}
\notag\\
&\qquad\qquad
+2\mathcal{M}^{(i)}_{bb}\mathcal{M}^{(j)}_{ca}\mathcal{M}^{(k)}_{ca}
+2\mathcal{M}^{(i)}_{cc}\mathcal{M}^{(j)}_{ab}\mathcal{M}^{(k)}_{ab}
+8\mathcal{M}^{(i)}_{bc}\mathcal{M}^{(j)}_{ca}\mathcal{M}^{(k)}_{ab}
\Big],
\\
Z_4
&=
\frac{3}{2}
\sum_{x,y,z\ \mathrm{pairwise\ distinct}}
\frac{
(\tau_{yyyz}-\tau_{zzzy})
\big[(B_z-B_x)\tau_{yyyz}+(B_x-B_y)\tau_{zzzy}\big]
}
{4(B_y-B_z)^2}.
\end{align}

Thus:
\begin{itemize}
\item $Z_1$, $Z_2$, and $Z_4$ are geometric and scale as $\omega^{-4}$;
\item $Z_3$ is cubic and scales as $\omega^{-3}$.
\end{itemize}

The cubic term splits as
\begin{equation}
Z_3 = Z_3^{(2)} + Z_3^{(3)},
\end{equation}
where
\begin{align}
\mathcal{K}_{abc}(p,q,r)
&=
\mathcal{M}^{(p)}_{aa}\mathcal{M}^{(q)}_{bb}\mathcal{M}^{(r)}_{cc}
+2\mathcal{M}^{(p)}_{aa}\mathcal{M}^{(q)}_{bc}\mathcal{M}^{(r)}_{bc}
\notag\\
&\quad
+2\mathcal{M}^{(p)}_{bb}\mathcal{M}^{(q)}_{ca}\mathcal{M}^{(r)}_{ca}
+2\mathcal{M}^{(p)}_{cc}\mathcal{M}^{(q)}_{ab}\mathcal{M}^{(r)}_{ab}
+8\mathcal{M}^{(p)}_{bc}\mathcal{M}^{(q)}_{ca}\mathcal{M}^{(r)}_{ab},
\end{align}
and
\begin{align}
Z_3^{(2)}
&=
\frac{1}{2}\sum_i
\frac{\Phi^{(3),\mathrm{red}}_{iii}}{\omega_i^3}
\mathcal{K}_{abc}(i,i,i)
\notag\\
&\quad
+\frac{1}{2}\sum_{i\neq j}
\frac{\Phi^{(3),\mathrm{red}}_{iij}}{\omega_i^2\omega_j}
\sum_{\pi\in\{(i,i,j),(i,j,i),(j,i,i)\}}
\mathcal{K}_{abc}(\pi_1,\pi_2,\pi_3),
\\
Z_3^{(3)}
&=
\frac{1}{2}\sum_{i<j<k}
\frac{\Phi^{(3),\mathrm{red}}_{ijk}}{\omega_i\omega_j\omega_k}
\sum_{\pi\in S_3}
\mathcal{K}_{abc}(\pi_1,\pi_2,\pi_3).
\end{align}

The formulas above are backend-faithful and make explicit that:
\begin{itemize}
\item the standard geometric sextic sector depends only on first derivatives of
$I^{-1}$ and harmonic rotational quantities, with overall harmonic scaling
$\omega^{-4}$;
\item the standard cubic sextic sector is linear in the reduced cubic constants
and scales overall as $\omega^{-3}$.
\end{itemize}
